\scanpage{197}\mspage{146}
complétant $X_*$, de telle façon que ce carré soit $W$-homotopiquement commutatif, i.e que $\varphi_Q=\alpha_Q\varphi_{Q_0}$ soit dans $W$. C'est d'autant plus plausible que dans $\mathrm{Cat}$, les sommes amalgamées ont \unsure{tendance} à ne pas avoir des bonnes propriétés homotopiques, même si l'une des flèches $\alpha_1,\alpha_2$ est mono, (voire une immersion ouverte), contrairement à ce qui se passe dans $\Delta^\wedge$, et plus généralement dans les catégories de la forme $Z^\wedge$ ($Z$ dans $\mathrm{Cat}$).

3.4. \uline{Critères pour les carrés homotopiquement cocartésiens}.

Les critères que j'ai dégagés sont tous contenus dans l'énoncé suivant :

\uline{\deleted{Proposition \unsure{11}}}\inserted{Théorème 4} Considérons un \inserted{diagramme} carré commutatif \unsure{dans|de} $\mathrm{Cat}$

% slanted
\margin{\uline{NB} Dans ce th., et les cor. 1, 2, 3, on n'a besoin que de la \uline{\unsure{forme} faible} de $W(7\text{bis})$, savoir que $\unsure{ir|i\circ r|i_0}\in W\Longrightarrow i\in W$ (cf p. \unsure{74|76}), (Loc($6'\mathrm a$) avec les notations de XVI)}

% slanted
\margin{\uline{NB} \unsure{Au total|En fait}, ce théorème \inserted{(\unsure{partie directe|partie droite})} est valable sans hypothèse $W(7\text{bis})$ ou $W(7\text{bis }\unsure{\text{faible}|\text{fort}})$ sur $W$, et il en est de $=$ de ses corollaires 1 à 3 (à l'exclusion des corollaires 4, 5), qui eux utilisent $W(7\text{bis})$ dans \unsure{tte} sa force).}

\diagram{\(Q={}\begin{tikzcd}[ampersand replacement=\&,row sep=1.6em,column sep=3em]
X_0\arrow[r,"\alpha_1"]\arrow[d,"\alpha_2"']\arrow[dr,dashed,"\unsure{f_0}"] \& X_1\arrow[d,"f_1"]\\
X_2\arrow[r,"f_2"'] \& Y
\end{tikzcd}\)}

et considérons le morphisme canonique

\diagram{\(\varphi_Q:\unsure{\mathscr X|\mathcal X}=\left(\int \begin{tikzcd}[ampersand replacement=\&,row sep=0pt,column sep=1.2em,cramped] \& X_1 \\ X_0\arrow[ur,"\alpha_1"]\arrow[dr,"\alpha_2"'] \& \\ \& X_2 \end{tikzcd}\right)\longrightarrow Y,\)}
% archive sheet number: 197
